\documentclass[11pt,a4paper]{article}

% --- Packages ---
\usepackage[margin=.75in]{geometry}
\usepackage[utf8]{inputenc}
\usepackage[T1]{fontenc}
\usepackage{lmodern} % Fix for font not found error
\usepackage{titlesec}
\usepackage{enumitem}
\usepackage{hyperref}
\usepackage{xcolor}
\usepackage{fancyhdr}
\usepackage{comment}
\usepackage{cleveref}
\usepackage[bottom]{footmisc}
\usepackage{graphicx}
\usepackage{longtable}
\usepackage{booktabs}
\usepackage{array}
\usepackage{calc}
\usepackage{etoolbox}

% Standard Pandoc setup
\providecommand{\tightlist}{%
  \setlength{\itemsep}{0pt}\setlength{\parskip}{0pt}}

\crefformat{section}{\S#2#1#3}
\crefformat{subsection}{\S#2#1#3}
\crefformat{subsubsection}{\S#2#1#3}
\crefformat{figure}{#2Figure~#1#3}
\crefformat{footnote}{#2\footnotemark[#1]#3}

\linespread{0.9} % Riduce lo spazio tra le linee

% --- Colors ---
\definecolor{primary}{RGB}{0, 0, 0}
\definecolor{links}{HTML}{0000EE}
\definecolor{blue}{HTML}{26428B}

\hypersetup{
    colorlinks=true,
    linkcolor=links,
    urlcolor=links,
    citecolor=links,
}

% --- Section Formatting ---
\titleformat{\section}
  {\normalfont\Large\bfseries\color{blue}}{\thesection}{1em}{}
\titleformat{\subsection}
{\normalfont\large\bfseries\color{blue}}{\thesubsection}{1em}{}
\titleformat{\subsubsection}
{\normalfont\bfseries\color{blue}}{\thesubsubsection}{1em}{}

% --- Header and Footer ---
\pagestyle{fancy}
\fancyhead[L]{\small Giuseppe Capasso}
\fancyhead[R]{\small devnet-workshop}
\fancyfoot[C]{\thepage}
\renewcommand{\headrulewidth}{0.4pt}

% --- Custom Title Command ---
\newcommand{\proposaltitle}[1]{
    \begin{center}
        \vspace*{0.1cm}
        {\LARGE \bfseries \color{blue} #1} \\[1.5ex]
        {\large \textbf{Giuseppe Capasso}} \\[1ex]
                \small{
            \vspace{0.1cm}
            \textbf{Materia:} devnet-workshop\\
        }
            \end{center}
    \vspace{0.3cm}
    \hrule height 0.5pt
    \vspace{0.5cm}
}

\begin{document}

\proposaltitle{REST API Lab Guide}

\section{REST API Lab Guide}\label{rest-api-lab-guide}

\subsection{Introduction}\label{introduction}

Welcome to the REST API Lab! In this lab, you will learn how to interact
with a REST API using Postman. Follow the instructions below to set up
Postman and complete the lab activities.

\subsection{Prerequisites}\label{prerequisites}

Before you begin, make sure you have the following: - Postman installed
on your computer - Basic knowledge of HTTP methods (GET, POST, DELETE,
PUT) - Basic understanding of REST APIs and JSON

\subsubsection{Installing Postman}\label{installing-postman}

Follow the instructions below to download and install Postman for your
specific platform:

\begin{itemize}
\tightlist
\item
  \textbf{Windows}:

  \begin{itemize}
  \tightlist
  \item
    Download Postman for Windows from
    \href{https://www.postman.com/downloads/}{here}.
  \item
    Run the installer and follow the installation instructions.
  \end{itemize}
\item
  \textbf{Mac}:

  \begin{itemize}
  \tightlist
  \item
    Download Postman for Mac from
    \href{https://www.postman.com/downloads/}{here}.
  \item
    Open the downloaded .dmg file and drag Postman to your Applications
    folder.
  \end{itemize}
\item
  \textbf{Linux (Debian-based distro)}:

  \begin{itemize}
  \item
    Download Postman for Linux from
    \href{https://www.postman.com/downloads/}{here}.
  \item
    Open a terminal and navigate to the directory where the downloaded
    file is located.
  \item
    Run the following commands:
  \end{itemize}

  \begin{itemize}
  \tightlist
  \item
    sudo tar -xzf Postman-linux-x64-8.0.9.tar.gz -C /opt sudo ln -s
    /opt/Postman/Postman /usr/bin/postman
  \end{itemize}
\end{itemize}

\subsection{Lab Activities}\label{lab-activities}

\subsubsection{Activity 1: Perform a GET
Request}\label{activity-1-perform-a-get-request}

\begin{enumerate}
\def\labelenumi{\arabic{enumi}.}
\tightlist
\item
  Open Postman.
\item
  In the request URL field, enter the endpoint URL:
  \texttt{http://example.com/devices}.
\item
  Select the HTTP method as \textbf{GET}.
\item
  Click the \textbf{Send} button to execute the request.
\item
  Review the response body to see the list of devices returned by the
  API.
\end{enumerate}

\subsubsection{Activity 2: Perform a GET Request for Device
Details}\label{activity-2-perform-a-get-request-for-device-details}

\begin{enumerate}
\def\labelenumi{\arabic{enumi}.}
\tightlist
\item
  Open Postman.
\item
  In the request URL field, enter the endpoint URL:
  \texttt{http://example.com/devices/\{id\}} (replace \texttt{\{id\}}
  with the ID of the device you want to retrieve details for).
\item
  Select the HTTP method as \textbf{GET}.
\item
  Click the \textbf{Send} button to execute the request.
\item
  Review the response body to see the details of the device.
\end{enumerate}

\subsubsection{Activity 3: Perform a DELETE Request (with
Authentication)}\label{activity-3-perform-a-delete-request-with-authentication}

\begin{enumerate}
\def\labelenumi{\arabic{enumi}.}
\tightlist
\item
  Open Postman.
\item
  In the request URL field, enter the endpoint URL:
  \texttt{http://example.com/devices/\{id\}} (replace \texttt{\{id\}}
  with the ID of the device you want to delete).
\item
  Select the HTTP method as \textbf{DELETE}.
\item
  Go to the \textbf{Authorization} tab.
\item
  Select \textbf{Bearer Token} from the Type dropdown.
\item
  Enter the static token \texttt{0InfNl5WZuR++sOUD4otAw==}.
\item
  Click the \textbf{Send} button to execute the request.
\item
  Verify that the device has been successfully deleted.
\end{enumerate}

\subsubsection{Activity 4: Perform a POST Request to Create a New Device
(with
Authentication)}\label{activity-4-perform-a-post-request-to-create-a-new-device-with-authentication}

\begin{enumerate}
\def\labelenumi{\arabic{enumi}.}
\tightlist
\item
  Open Postman.
\item
  In the request URL field, enter the endpoint URL:
  \texttt{http://example.com/devices}.
\item
  Select the HTTP method as \textbf{POST}.
\item
  Go to the \textbf{Authorization} tab.
\item
  Select \textbf{Bearer Token} from the Type dropdown.
\item
  Enter the static token \texttt{0InfNl5WZuR++sOUD4otAw==}.
\item
  Go to the \textbf{Body} tab.
\item
  Select \textbf{Raw} and choose \textbf{JSON} from the dropdown.
\item
  Enter the JSON object representing the new device:
\end{enumerate}

\begin{verbatim}
{
    "serial": "CAT-2960-5483221",
    "model": "Catalyst-2960x-48p",
    "type": "SWL2",
    "interfaces": [
        {
            "name": "VLAN1",
            "address": "192.168.1.100",
            "netmask": "255.255.255.0",
            "description": "A VLAN",
            "status": "Administrative down"
        }
    ]
}
\end{verbatim}

\subsubsection{Activity 5: Perform a PUT Request to Update a
Device}\label{activity-5-perform-a-put-request-to-update-a-device}

\begin{enumerate}
\def\labelenumi{\arabic{enumi}.}
\tightlist
\item
  Open Postman.
\item
  In the request URL field, enter the endpoint URL:
  \texttt{http://example.com/devices/\{id\}} (replace \texttt{\{id\}}
  with the ID of the device you want to update).
\item
  Select the HTTP method as \textbf{PUT}.
\item
  Go to the \textbf{Authorization} tab.
\item
  Select \textbf{Bearer Token} from the Type dropdown.
\item
  Enter the static token \texttt{0InfNl5WZuR++sOUD4otAw==}.
\item
  Go to the \textbf{Body} tab.
\item
  Select \textbf{Raw} and choose \textbf{JSON} from the dropdown.
\item
  Enter the updated JSON object representing the device:
\end{enumerate}

\begin{verbatim}
{
    "serial": "UPDATED_SERIAL_NUMBER",
    "model": "UPDATED_MODEL",
    "type": "UPDATED_TYPE",
    "interfaces": [
        {
            "name": "VLAN1",
            "address": "UPDATED_IP_ADDRESS",
            "netmask": "UPDATED_NETMASK",
            "description": "UPDATED_DESCRIPTION",
            "status": "UPDATED_STATUS"
        }
    ]
}
\end{verbatim}

\subsection{Conclusion}\label{conclusion}

Congratulations! You have completed the REST API lab using Postman. You
have learned how to perform GET, PUT, DELETE, and POST requests, as well
as how to handle authentication using Bearer tokens.

If you have any questions or need assistance, feel free to ask your
instructor or classmates.

Happy API testing!

\end{document}
